\documentclass[11pt]{article}
\usepackage[margin=1in]{geometry}
\usepackage[T1]{fontenc}
\usepackage[utf8]{inputenc}
\usepackage{lmodern}
\usepackage{microtype}
\usepackage{amsmath,amssymb,amsthm,mathtools}
\usepackage{booktabs}
\usepackage{array}
\usepackage{longtable}
\usepackage{graphicx}
\usepackage{xcolor}
\usepackage{hyperref}
\usepackage[nameinlink,noabbrev]{cleveref}
\usepackage{enumitem}
\usepackage{listings}
\usepackage{csquotes}

\hypersetup{
    colorlinks=true,
    linkcolor=blue,
    citecolor=blue,
    urlcolor=blue
}

\setlist[itemize]{leftmargin=1.5em}
\setlist[enumerate]{leftmargin=1.75em}

\newcommand{\trellis}{\textsc{Trellis}}
\newcommand{\platformrequest}{\texttt{PlatformRequest}}
\newcommand{\semanticcontract}{\texttt{SemanticContract}}
\newcommand{\valuationcontext}{\texttt{ValuationContext}}
\newcommand{\productir}{\texttt{ProductIR}}
\newcommand{\eventprogramir}{\texttt{EventProgramIR}}
\newcommand{\controlprogramir}{\texttt{ControlProgramIR}}
\newcommand{\familyir}{family IR}
\newcommand{\code}[1]{\texttt{#1}}


\title{Trellis II: Mathematical and Computational Lanes for Semantic Derivatives Pricing}
\author{Draft}
\date{\today}

\begin{document}
\maketitle

\begin{abstract}
This paper should explain the mathematical and computational substrate that
lets one semantic request lower onto multiple deterministic pricing families.
The main claim is that \trellis{} succeeds by combining contract algebra,
event/control compilation, bounded family IRs, and reusable numerical kernels
rather than by relying on a single universal solver representation.
\end{abstract}

\tableofcontents

\section{Introduction}

Frame the paper as the bridge between rich product semantics and bounded
numerical execution.

\section{Why Bounded Family IRs}

Explain why the system does not use one monolithic universal pricing IR.

\section{Contract Algebra and Analytical Support}

Cover:
\begin{itemize}
    \item contracts as syntax
    \item valuation as semantics
    \item semantics-preserving rewrites
    \item reusable analytical kernels
\end{itemize}

\section{Event and Control as Numerical Interface}

Cover:
\begin{itemize}
    \item \eventprogramir{}
    \item \controlprogramir{}
    \item how shared semantic programs project into family-specific forms
\end{itemize}

\section{Family Lanes}

Suggested subsections:
\begin{itemize}
    \item analytical
    \item exercise lattice
    \item PDE
    \item event-aware Monte Carlo
    \item transforms
    \item basket credit / copula
\end{itemize}

\section{Calibration and Market Binding}

Explain how model binding, quote regimes, calibration surfaces, and market
context stay explicit rather than route-local.

\section{Case Studies}

Suggested cases:
\begin{itemize}
    \item transform family hardening
    \item event-aware Monte Carlo swaption path
    \item short-rate helper extraction
\end{itemize}

\section{Support Boundary and Limitations}

Use \code{LIMITATIONS.md} directly to keep the paper honest about current
coverage, especially for calibration, risk, xVA, and performance.

\section{Conclusion}

Restate the main result: the computational novelty lies in family-specific
semantic lowering onto deterministic quantitative kernels.

\end{document}
